\subsection{Tái lập thực nghiệm}
\subsubsection{Mã nguồn và môi trường}
Notebook công khai kèm mã và cấu hình:
\begin{center}
{\small\url{https://www.kaggle.com/code/kieuthithutrang/gtsrb-low-light-enhancement-for-traffic-signs}}
\end{center}
Nhóm tải dữ liệu từ archive GTSRB chính thức. Mã và trọng số pretrained của các mô hình tăng cường được ghi nguồn, commit và SHA-256 để xác định phiên bản. Trong các phép so sánh, nhóm giữ nguyên cách chia tập, mã ảnh, mức suy giảm, noise seed và ba checkpoint. Tham số xử lý được chọn trên validation trước khi chạy test. Người đọc có thể xem hướng dẫn chạy trong \path{README.md} và các cấu hình trong \path{configs/}.

\subsubsection{Bằng chứng và tính nhất quán}
Gói kết quả lưu giao thức, trạng thái chạy, dự đoán CSV và ma trận nhầm lẫn NPZ. Nhóm đã tính lại Accuracy, Macro-F1 và Weighted-F1 từ 5.847.990 dòng dự đoán thuộc bốn phạm vi. Sai số so với bảng tổng hợp không quá $10^{-8}$, và hash ba CNN khớp giữa các phạm vi. Những đối chứng dùng lại đều được ghi nguồn, không xem là quan sát độc lập mới. Việc đối chiếu nhãn và group ID giữa các lượt giúp kiểm tra tính nhất quán, nhưng không thay thế kiểm tra nhãn trên ảnh nguồn.

Nhóm lưu bằng chứng và thông tin nguồn trong \path{report/evidence/final/}. ZIP kết quả không chứa ảnh GTSRB, cache hoặc trọng số pretrained của các mô hình tăng cường. Người chạy lại cần lấy những đầu vào đó theo nguồn và cấu hình đã ghi. ZIP vẫn giữ các checkpoint CNN do nhóm huấn luyện. SHA-256 giúp phát hiện file có thay đổi, nhưng tự nó không chứng minh kết luận nghiên cứu là đúng.
